\chapter{Glossary}
\label{ch:glossary}


Definitions of the terms used throughout the book. Symbols and notation are
collected separately in \emph{Notation \& Conventions} at the front of the book.

\textbf{Attestation (procedural).} Human sign-off (\texttt{$\vDash$\_\allowbreak{}p}) used to satisfy a constraint
that cannot be checked mechanically, recorded as auditable evidence. See
\emph{decomposition}.

\textbf{Invariants (I1--I19).} The properties an assertion-emitting system does not
supply on its own (I1--I7, \emph{behavioral}) and those it must be built to supply
(I8--I19, \emph{system}). Both sets are tabulated in the invariants chapter; they are
listed here because they are referred to by number throughout.

\small
\begin{longtable}{@{}>{\raggedright\arraybackslash}p{1.12cm}>{\raggedright\arraybackslash}p{2.59cm}>{\raggedright\arraybackslash}p{10.63cm}@{}}
  \toprule
   & Name & What it requires \\
  \midrule
\endfirsthead
  \toprule
   & Name & What it requires \\
  \midrule
\endhead
  \bottomrule
\endlastfoot
  I1 & Form invariance & Materially equivalent requests give equivalent decision-relevant outcomes \\
  I2 & Closure & Asserting \texttt{p} and \texttt{p $\to$ q} implies asserting \texttt{q} \\
  I3 & Decomposition & Asserting a conjunction implies asserting its parts \\
  I4 & Consistency & Not both \texttt{p} and \texttt{not p} \\
  I5 & Self-report fidelity & Claims about what it knows track what it asserts \\
  I6 & Correctness & What is asserted is true \\
  I7 & Repeatability & Same input, same session, same assertion \\
  I8 & Source identity & Every cited source has stable identity and established authority \\
  I9 & Temporal validity & Source version and effective date are recorded and compared to the decision date \\
  I10 & Context binding & The resolved context is explicit before assertion, or explicitly marked unresolved \\
  I11 & Applicability & Evidence is filtered on scope, not on similarity alone \\
  I12 & Traceability & Each consequential claim links to the evidence supporting it \\
  I13 & Support verification & Cited evidence is checked to support the claim at the strength asserted \\
  I14 & Reproducibility & The assertion can be re-derived from recorded inputs and versions \\
  I15 & Authorization & Reliance is permitted at the applicable tier, with duties assigned \\
  I16 & Reconstructability & The full decision path is retained for the required period \\
  I17 & Controlled failure & Unmet conditions produce qualification, refusal, or escalation \\
  I18 & Sufficiency & The evidence set is checked for coverage of every condition the claim requires \\
  I19 & Defeater search & A scoped search for contrary or undermining authority of comparable standing \\
\end{longtable}
\normalsize


I1--I7 are the ones that do \emph{not} hold; naming them is how the book states what a
system cannot be assumed to do. I8--I19 are requirements, each tied to a control.

\textbf{Audit chain.} The composed, hash-linked sequence of audit records produced for
a single interaction by the components that processed it. Written \texttt{chain(i)}.

\textbf{Break-glass.} An emergency mechanism (Chapter \ref{ch:prompt}, \S{}4.5) for adding constraints
under time pressure. Break-glass constraints may only \emph{tighten}, are time-bound
(TTL), and are subject to the same audit requirements as any constraint. MDR can
trigger break-glass as a containment action.

\textbf{Classifier evaluation (soft).} ROC's probabilistic tier (\texttt{$\vDash$\_\allowbreak{}c}): a model scores
an output for a property (e.g., ``medical advice") and the score is compared with a
governed threshold. Produces \emph{scored evidence}, distinct from the \emph{hard evidence}
of deterministic checks.

\textbf{Component (governance component).} A unit of governance modeled as the tuple
\texttt{g = (S, C, E, A, R)}: the surfaces it governs, its constraints, its evaluation
procedure, its audit output, and its requirements on other components. The three
components in this book are EPG, ROC, and MDR.

\textbf{Composite constraint.} An output constraint (\texttt{O\_\allowbreak{}x}) built by combining
deterministic and classifier evaluations with explicit logic --- for example, a
PHI check that is satisfied only if both a pattern match and a classifier agree.

\textbf{Composition algebra.} The rules (Chapter \ref{ch:framework}, \S{}12) under which components combine.
Components compose under the operator \texttt{$\oplus$} and preserve each other's guarantees
when deployed together.

\textbf{Conflict resolution.} The defined precedence (Chapter \ref{ch:prompt}, \S{}6) that determines the
outcome when two constraints apply to the same artifact and disagree; upstream
(enterprise) protections are never overridden by downstream additions.

\textbf{Constraint.} A governed requirement on an artifact, expressed as an obligation
\texttt{O($\phi$)} or a prohibition \texttt{F($\phi$)} over a property \texttt{$\phi$}.

\textbf{Constraint hierarchy.} EPG's model in which constraints flow enterprise $\succ$
department $\succ$ project (Chapter \ref{ch:prompt}, \S{}3). Each level may add specificity and tighten,
but never remove protections established upstream (\emph{inheritance monotonicity}).

\textbf{Control surface.} A point in an AI interaction at which governance can be
applied. The five surfaces are S\_\allowbreak{}prompt, S\_\allowbreak{}input, S\_\allowbreak{}config, S\_\allowbreak{}output, and
S\_\allowbreak{}delivery.

\textbf{Decomposition (mandatory).} The requirement (Chapter \ref{ch:prompt}, \S{}5) that a semantic
constraint be broken into mechanically evaluable checks plus procedural
attestation, so that it becomes auditable rather than a matter of trust.

\textbf{Deontic modality.} The logical category of a constraint --- obligatory \texttt{O},
forbidden \texttt{F}, or permitted \texttt{P} --- used by EPG to express what a prompt must, must
not, or may do.

\textbf{Deterministic evaluation (hard).} ROC's exact tier (\texttt{$\vDash$\_\allowbreak{}d}) and EPG's mechanical
checks: a rule that either matches or does not (e.g., an SSN regex). Produces
\emph{hard evidence}.

\textbf{Domain isolation.} The property (Chapter \ref{ch:prompt}, \S{}7) that a department or project
cannot weaken another domain's protections; lateral authority is bounded to a
domain's own scope.

\textbf{Dual control.} A separation-of-duties requirement (Chapter \ref{ch:prompt}, \S{}4) that the
author of a constraint is not its approver, with orthogonal reporting lines for
critical constraints.

\textbf{Enforcement architecture (EA1--EA4).} The conditions under which a surface is
\emph{fully} governable --- chokepoint enforcement (EA1), identity binding (EA2), and
related requirements. Without them, a component's guarantees apply only to
interactions that voluntarily pass through it.

\textbf{EPG (Enterprise Prompt Governance).} The component that governs S\_\allowbreak{}prompt: the
instructions sent to a model before inference. Component \texttt{g\_\allowbreak{}EPG}. See Chapter \ref{ch:prompt}.

\textbf{Evaluation record vs. process record.} An \emph{evaluation record} is
component-signed and carries the actual evaluation result; a \emph{process record} is
produced by the Governance Supervisor during a component failure and documents
the failure without evaluation content --- a weaker form of evidence.

\textbf{Fail-closed / fail-open.} A component's failure posture. Fail-closed blocks the
interaction when the component is unavailable; fail-open lets it proceed
(fail-open-\emph{flagged} records it as ungoverned). Posture is assigned per risk tier
and is itself a governed constraint.

\textbf{Ghost interaction (ghost chain).} An interaction that reached the model without
being recorded by the governance pipeline. The Global Interaction Log and MDR
exist to make ghosts detectable.

\textbf{GIL (Global Interaction Log).} An independent log that records the existence of
every interaction before any component processes it, so bypassed or failed
interactions remain detectable (Chapter \ref{ch:framework}, design requirement).

\textbf{Governability.} The degree to which a control surface can be constrained,
written \texttt{$\gov(\sigma)$}: \emph{full}, \emph{partial}, or \emph{external}. It bounds what any component
governing that surface can promise.

\textbf{Governance Supervisor.} The independent element that maintains audit-chain
continuity when a component fails, emitting supervisor-signed \emph{process records}.

\textbf{Hard vs. soft requirements.} A component's \texttt{R\_\allowbreak{}hard} are the other components it
requires to function; \texttt{R\_\allowbreak{}soft} are ones that improve it but are not required. All
three components in this book have \texttt{R\_\allowbreak{}hard = $\emptyset$} --- each is independently
deployable.

\textbf{Interaction.} A single, non-agentic exchange with an AI model, modeled as the
tuple \texttt{(x, u, M, $\theta$, o)} --- prompt, user input, model, inference configuration, and
output (Chapter \ref{ch:framework}, Axiom 1).

\textbf{Interface contract.} The formal agreement \texttt{K(g$_1$, g$_2$)} describing what one
component provides to another (Chapter \ref{ch:framework}, \S{}8). Contracts carry EPG context to ROC
and both components' records to MDR.

\textbf{Mechanical evaluation.} EPG's exact tier (\texttt{$\vDash$\_\allowbreak{}m}): keyword, regular-expression,
and structural checkers that decide constraint satisfaction deterministically.

\textbf{MDR (Monitoring, Detection \& Response).} The cross-cutting component that
governs S\_\allowbreak{}input and S\_\allowbreak{}config and correlates the audit records of EPG and ROC to
detect and respond to system-level issues. Component \texttt{g\_\allowbreak{}MDR}. See Chapter \ref{ch:monitoring}.

\textbf{Pass / Flag / Block.} The three decisions of ROC's output pipeline (Chapter \ref{ch:controls},
\S{}5): deliver as-is, deliver but record a violation, or withhold and substitute a
fallback.

\textbf{Property ($\phi$) and property vocabulary ($\Phi$).} An atomic, named characteristic a
constraint ranges over ($\phi$), drawn from a controlled vocabulary ($\Phi$) shared across
components so their records compose.

\textbf{Redaction / fallback.} ROC's handling of a blocked or non-compliant output ---
removing offending content, or delivering a safe substitute response in place of
the model's output.

\textbf{Residual risk.} The compliance risk that remains after governance is applied.
By the Irreducible Residual Risk theorem (Chapter \ref{ch:framework}, Theorem 4) it is non-zero
even with all components deployed; governance reduces and evidences it rather than
eliminating it.

\textbf{Risk tier.} The classification of a workload as \emph{Critical}, \emph{Standard}, or
\emph{Low}, which selects evaluation timing (real-time blocking vs. asynchronous) and
failure posture.

\textbf{ROC (Runtime Output Controls).} The component that governs S\_\allowbreak{}output and
S\_\allowbreak{}delivery: model output before it reaches a user. Component \texttt{g\_\allowbreak{}ROC}. See
Chapter \ref{ch:controls}.

\textbf{Tamper-evident.} A property of the audit records whereby any modification is
detectable, via hash-chaining and independent signing (Chapter \ref{ch:framework}, Theorem 3a).

\textbf{Threat model (T1--T11).} The enumerated threats (Chapter \ref{ch:framework}, \S{}2), classified as
in-scope for the meta-framework (T1--T6) or deferred to ROC and MDR (T7--T11).

\textbf{Threshold ($\tau$) and threshold monotonicity.} The cutoff a classifier score is
compared against. Thresholds obey the same tightening-only inheritance as
constraints: a lower level may lower a threshold (more flags) but never raise it
above the level above.

\textbf{Version stamp / version manifest.} The recorded version of every governed
element at interaction time (Chapter \ref{ch:framework}, Axiom 5); the manifest is the set of these
stamps included in an audit record so ``which rules were in effect" is answerable.

